\section{Discussion}\label{sec:discussion}

\subsection{Architectural and RI Implications}
The main architectural lesson is that the gateway should be treated as a semantic boundary, not as a transparent relay. Constrained devices can remain focused on protected communication, local execution, and evidence emission, while the gateway handles identity validation, protocol translation, and status propagation toward Kubernetes. This separation avoids pushing cloud-facing semantics into firmware and gives the control plane a stable interface even if the southbound protocol evolves.

The evaluation confirms that protocol success alone is not sufficient: enrollment, liveness, and deployment evidence must be interpreted consistently before they become trustworthy control-plane state. This supports the central design decision to separate desired state, gateway-derived evidence, and controller reconciliation. Future extensions such as multi-gateway placement, staged rollout, or retry policies should preserve this separation rather than allowing multiple components to write conflicting interpretations of device status.

For sustainable digital research infrastructures, the same separation creates explicit attachment points for energy and carbon policies. Placement decisions can be expressed as desired state, gateway and device observations can be treated as evidence, and sustainability controllers can consume the same status model without requiring constrained devices to run additional cloud-native agents. The measured baselines in Section~\ref{sec:sustainability-metrics}---byte-scale southbound traffic, sub-20~MiB gateway mediation, and no kubelet on the edge---suggest that subsequent green-experimentation studies can focus on policy and instrumentation rather than on redesigning the orchestration substrate.

\subsection{Validity and Scope}
The validation targets system-level orchestration behavior rather than real-time performance, radio behavior, or sustainability optimization. Emulation and scripted trials are appropriate for exercising enrollment, deployment, status convergence, and reconciliation semantics, while physical boards, wireless links, and power meters require additional characterization. The results should therefore be interpreted as evidence that the orchestration model is coherent under the evaluated workflow.

External validity is bounded by the single-node K3s deployment, one gateway path, and a representative minimal workload. Larger deployments introduce gateway contention, distributed failure modes, network partitions, and device mobility. A supplementary scaling attempt with three Renode machines on the shared TAP test bed kept only one TLS session active in the gateway runtime view; fleet-scale energy characterization therefore requires isolated network segments or physical boards. Construct validity is also scoped: the study measures secure attachment, observable liveness, deployment completion, latency, wire sizes, firmware footprint, and control-plane RAM/CPU samples, while direct energy consumption and long-duration reliability remain outside the current evaluation.

Conclusion validity is strengthened by the transactional success criterion: enrollment, heartbeat visibility, and deployment must succeed within the same trial. This avoids overstating correctness from isolated protocol stages. Broader claims about optimality, scalability, or comparative advantage require additional experiments against alternative orchestration designs and physical-device deployments.

\subsection{Reproducibility}
The evaluation is designed to be repeatable on a commodity Linux host with Docker and K3s. Renode makes the embedded side deterministic enough for regression-style validation, although it remains complementary to future physical-board experiments. Evidence traceability is maintained across multiple layers: gateway logs, API-visible resource states, and Kubernetes CRD status describe the same workflow from different perspectives, so that protocol success and control-plane convergence are not inferred from a single eventually consistent field.
